% =============================================================
% Introduction générale (~3 pages) — sans numérotation de section.
% Cinq mouvements imposés (note pédagogique ESPRIT) : sujet,
% problématique (sans évoquer aucun résultat), démarche, objectifs,
% annonce du plan. S'y ajoutent, depuis la refonte du 29/08/2026 :
% l'encadré de convention de désignation, l'annonce du fil rouge et
% la carte du mémoire (figure 0).
% La problématique est rigoureusement identique à celle de la
% section 1.4.1 : toute modification doit être répercutée aux deux
% endroits.
% =============================================================
\chapter*{Introduction générale}
\addcontentsline{toc}{chapter}{Introduction générale}
\markboth{Introduction générale}{Introduction générale}

\textbf{Sujet.} La sécurité des systèmes d'information repose traditionnellement sur un modèle
périmétrique : filtrer en entrée, faire confiance à l'intérieur. Ce modèle a fait ses preuves
quand les ressources étaient centralisées et les utilisateurs sur le réseau local. Il ne les
fait plus aujourd'hui. Les applications migrent dans le cloud, les données sont accessibles de
n'importe où, et les attaques les plus coûteuses exploitent un accès légitime détourné plutôt
qu'une intrusion frontale.

MENAL-SARL, entreprise de services informatiques implantée en Mauritanie et en France, a
identifié ce risque sur l'une de ses applications clientes : ELSON, une plateforme de collecte
participative pour les langues nationales de Mauritanie. Cette application collecte des données
d'identité officielle et des enregistrements vocaux, qui constituent des données biométriques.
Elle fonctionne, mais elle repose sur un serveur unique, avec confiance implicite entre tous ses
composants et sans mécanisme de reproductibilité.

\begin{keybox}[title={Convention de désignation}]
Les composants du socle sont désignés par leur \emph{capacité} --- ce qu'ils font --- plutôt que
par leur nom commercial, afin que l'architecture reste lisible indépendamment du fournisseur. Le
produit qui réalise chaque capacité est indiqué entre parenthèses à sa première occurrence, et la
correspondance complète figure au tableau d'inventaire technologique de
\cref{sec:inventaire}. Les
familles d'identifiants employées dans le corps sont récapitulées dans la page
\emph{Conventions de notation}, en tête de document.
\end{keybox}

\textbf{Problématique.} Comment concevoir, déployer et surtout \emph{prouver} un socle
d'hébergement cloud aligné sur les principes Zero Trust et sur une chaîne de livraison
DevSecOps ? La question se pose dans les conditions réelles d'une entreprise émergente : une
seule personne pour l'exploitation, un budget de quelques dizaines d'euros par mois, aucune
solution de supervision commerciale. Et quel prix faut-il payer, en faux positifs, en latence de
détection et en complexité, pour atteindre ce résultat ?

La question n'est pas seulement technique. Elle comporte un coût : une architecture qui ne
mesure pas ce qu'elle coûte n'est pas une architecture d'ingénieur. Elle comporte aussi une
exigence de preuve. Chaque contrôle doit être associé à un test qui établit son efficacité,
plutôt qu'à une simple déclaration de conformité.

\textbf{Démarche.} Le projet suit une démarche en huit phases, PH0 à PH7, chacune fermée par un
critère de sortie vérifiable et détaillée en \cref{subsec:demarche}. Un audit technique a d'abord
établi l'état réel de l'existant. L'analyse de risque a ensuite identifié les scénarios
d'attaque les plus critiques, traduits en vingt exigences vérifiables. L'architecture a été
conçue autour de quatre principes directeurs : le principe de justification (PR1), qui impose à
tout composant de justifier son existence ; le principe d'identité (PR2), qui fait de l'identité
le plan de contrôle ; le principe d'intégrité de la preuve (PR3) ; et le principe de description
en code (PR4). La réalisation a construit le socle en infrastructure décrite en code, avec deux
chaînes de livraison et une chaîne de détection enrichie par un modèle de représentation
sémantique. Une campagne de validation confronte enfin les résultats aux objectifs mesurables
fixés au départ.

\textbf{Objectifs.} Six objectifs mesurables structurent le projet : supprimer la confiance
implicite du réseau (O1), rendre l'infrastructure reproductible (O2), contrôler la chaîne de
livraison (O3), détecter les événements de sécurité (O4), enrichir les alertes de façon
déterministe (O5), rester exploitable et économiquement soutenable (O6). Chaque objectif est
associé à un indicateur, une cible et un moyen de vérification. Le \cref{ch:validation} y répond point par
point.

\textbf{Fil conducteur.} Une même attaque traverse ce mémoire de bout en bout. Le 19/08/2026,
treize requêtes malveillantes ont été envoyées contre l'application hébergée en environnement de
recette. Le lecteur les rencontrera successivement sous quatre formes : comme scénario de menace
dérivé de l'analyse de risque au \cref{ch:besoins-menaces}, comme séquence de refus au périmètre au \cref{ch:conception},
comme règle de détection au \cref{ch:realisation}, et comme mesure au \cref{ch:validation}. Suivre ce cas unique d'un
chapitre à l'autre est la façon la plus courte de vérifier que la conception, la réalisation et
la mesure parlent bien du même système.

\textbf{Plan du mémoire.} Le \cref{ch:cadre-existant} établit le contexte et la problématique sur des faits
vérifiés, et présente l'entreprise d'accueil ainsi que l'application qui sert de cas réel. Le
\cref{ch:etat-art} examine l'état de l'art, justifie les choix technologiques et publie l'inventaire des
briques structurantes du socle avec, pour chacune, le prix payé. Le \cref{ch:besoins-menaces} modélise les
menaces à partir d'un diagramme de flux de données et de ses frontières de confiance, et en
déduit les vingt exigences de sécurité. Le \cref{ch:conception} présente l'architecture Zero Trust qui les
réalise, vue par vue et couche par couche. Le \cref{ch:realisation} rend compte de la réalisation :
infrastructure en code, chaîne de livraison applicative, chaîne de provisionnement de
l'infrastructure, chaîne de détection, tableau de bord et API de supervision, accueil d'un
locataire. Le \cref{ch:validation} confronte l'ensemble à l'épreuve de la mesure --- tests adverses,
évaluation de la détection, performance, résilience et coût --- puis énonce les conditions
d'industrialisation du socle en service.

La \cref{fig:carte-memoire} donne la carte de cette progression. Elle se lit de haut en bas :
chaque chapitre consomme le livrable du précédent et en produit un nouveau, de sorte qu'aucun
maillon ne repose sur une affirmation non établie en amont. C'est cette chaîne, et non
l'accumulation de contrôles techniques, qui fait la démonstration.

\textbf{Sur le volume de ce document.} Ce mémoire est plus long que le format habituel, et
c'est un choix assumé plutôt qu'un débordement. La thèse défendue n'est pas qu'un socle a été
construit, mais que chaque contrôle qui le compose est adossé à l'épreuve qui en établit
l'efficacité : une chaîne qui va d'une menace à une exigence, de l'exigence à une décision
d'architecture, de la décision au code qui la réalise, et du code au protocole qui la met à
l'épreuve. Une telle chaîne ne se démontre pas par échantillonnage --- il suffit d'un maillon
manquant pour qu'elle cesse d'en être une. C'est pourquoi les vingt exigences, les dix-sept
décisions d'architecture, les écarts d'implémentation et les vingt protocoles sont publiés au
complet, y compris lorsqu'ils concluent contre le socle.
Le lecteur qui souhaite aller vite peut s'en tenir aux registres, qui donnent l'état du système
sous forme tabulaire : la matrice de traçabilité du \cref{ch:besoins-menaces}, le registre des
décisions et celui des écarts du \cref{ch:conception}, et le tableau de synthèse de la campagne
au \cref{ch:validation}. Le corps du texte en est le développement argumenté, et les annexes le
matériel de reproduction --- protocoles détaillés, extraits de configuration et jeu
d'évaluation --- tenu hors du fil de lecture précisément pour ne pas l'alourdir.

\begin{figure}[!htbp]
\centering
\begin{tikzpicture}[
  font=\footnotesize,
  bloc/.style={nsysteme, text width=4.05cm, align=left},
  prod/.style={text width=8.35cm, align=left, font=\scriptsize,
               text=menalencre, anchor=west}]

  \node[bloc] (c1) at (0,0.00) {\textbf{Chapitre~1}\\Contexte et existant};
  \node[prod] at (2.35,0.00) {\textbf{Produit :} un audit daté sur 41~critères, cinq carences C1--C5, six objectifs mesurables O1--O6.};
  \node[bloc] (c2) at (0,-1.32) {\textbf{Chapitre~2}\\État de l'art et choix};
  \node[prod] at (2.35,-1.32) {\textbf{Produit :} quatre verrous V1--V4, l'inventaire des briques structurantes et le prix payé pour chacune.};
  \node[bloc] (c3) at (0,-2.64) {\textbf{Chapitre~3}\\Besoins et menaces};
  \node[prod] at (2.35,-2.64) {\textbf{Produit :} un diagramme de flux à cinq frontières de confiance, vingt exigences EX1--EX20 et la matrice de traçabilité.};
  \node[bloc] (c4) at (0,-3.96) {\textbf{Chapitre~4}\\Conception};
  \node[prod] at (2.35,-3.96) {\textbf{Produit :} quatre principes PR1--PR4, sept flux F1--F7, le registre des décisions D00--D16 et celui des écarts.};
  \node[bloc] (c5) at (0,-5.28) {\textbf{Chapitre~5}\\Réalisation};
  \node[prod] at (2.35,-5.28) {\textbf{Produit :} le socle décrit en code, deux chaînes de livraison, sept règles de détection R1--R7, les deux logiciels écrits.};
  \node[bloc] (c6) at (0,-6.60) {\textbf{Chapitre~6}\\Validation et mesure};
  \node[prod] at (2.35,-6.60) {\textbf{Produit :} vingt résultats de tests T1--T20, les mesures de détection, de résilience et de coût, les conditions d'industrialisation.};

  \draw[fluxdonnee] (c1.south) -- (c2.north);
  \draw[fluxdonnee] (c2.south) -- (c3.north);
  \draw[fluxdonnee] (c3.south) -- (c4.north);
  \draw[fluxdonnee] (c4.south) -- (c5.north);
  \draw[fluxdonnee] (c5.south) -- (c6.north);

\end{tikzpicture}
\caption[Carte du mémoire : chaque chapitre consomme le livrable du précédent]{Carte du mémoire : chaque chapitre consomme le livrable du précédent. Aucun
maillon de la démonstration ne repose sur une affirmation qui ne soit établie en amont.}
\label{fig:carte-memoire}
\end{figure}
